\input{../tex-inputs/latex-leseansicht-vorspann}

\section[Arthur Schnitzler an Gustav Schwarzkopf, 18. 6. 1917]{L04401 Arthur Schnitzler an Gustav Schwarzkopf, 18. 6. 1917}
\nopagebreak\mylabel{L04401v}
\rehead{ }\normalsize\beginnumbering\briefempfaengerindex{Schwarzkopf, Gustav@\textsc{Schwarzkopf, Gustav}!zzzSchnitzler, Arthur@\emph{von Arthur Schnitzler}!1917-06-181@{18. 6. 1917}|(be}
\toendnotes[C]{\smallbreak\pagebreak[2]}
\correspDesc{Versand  durch Arthur Schnitzler am 18. 6. 1917 in Böckstein
\newline{}Übermittlung  in Böckstein
\newline{}Erhalt  durch Gustav Schwarzkopf im Zeitraum [19. 6. 1917 – 23. 6. 1917?] in Wien}\toendnotes[C]{\smallbreak}
\Standort{DLA, A:Schnitzler, HS.1985.1.1897.}
\physDesc{Bildpostkarte, 216 Zeichen
\newline{}Handschrift: Bleistift, deutsche Kurrent
\newline{}Versand: 1) Stempel: »\nobreak{}\oindex{Valeriehaus@\textbf{Valeriehaus}|pwk}Valeriehaus Nassfeld\nobreak{}«.   2) Stempel: »\nobreak{}\oindex{Böckstein@\textbf{Böckstein}|pwk}B\textcolor{gray}{ö}ckstein\nobreak{}«. }\toendnotes[C]{\smallbreak}\pstart{}{\pb}Herrn Gustav Schwarzkopf\pend{}\pstart{}Wien I\oindex{I., Innere Stadt@\textbf{I., Innere Stadt}|pw}\pend{}\pstart{}Tiefer Graben 17\oindex{Wien@\textbf{Wien}!I., Innere Stadt@\textbf{I., Innere Stadt}!Tiefer Graben 17@\textbf{Tiefer Graben 17}|pw}\pend{}{\bigskip}
\pstart
           \centering{}{\pb}\textcolor{gray}{\textbf{NASSFELD\oindex{Nassfeld@\textbf{Nassfeld}|pw} mit Keesause\oindex{Keesause@\textbf{Keesause}|pw}. Schlapperebenspitze\oindex{Schlapperebenspitzen@\textbf{Schlapperebenspitzen}|pw} u Muraukopf\oindex{Murauer Köpfe@\textbf{Murauer Köpfe}|pw}.}}\pend
           \vspace{1em}
\pstart
           \raggedleft{}{\pb}18. 6. 17\pend
           \vspace{0.5em}
\pstart
           Herzliche Grüße und auf baldgs Wiederſehen. Es wird  immer{ }ſchöner hier, aber \label{K_L04401-1v}\edtext{Ende
               der Woche gehts doch nach Wien\oindex{Wien@\textbf{Wien}|pw}}{\lemma{\textnormal{\emph{Ende … Wien}}}\Cendnote{\textnormal{Er traf am 22. 6. 1917 wieder
               in Wien\oindex{Wien@\textbf{Wien}|pwk} ein.}}}\label{K_L04401-1}.\pend
           
\pstart
           Von einer kurzen einſamen Bergwanderung.\pend
           \pstart Ihr \spacefill\mbox{Arthur}\pend{}\selectlanguage{ngerman}\endnumbering\briefempfaengerindex{Schwarzkopf, Gustav@\textsc{Schwarzkopf, Gustav}!zzzSchnitzler, Arthur@\emph{von Arthur Schnitzler}!1917-06-181@{18. 6. 1917}|)be}\mylabel{L04401h}
\begin{anhang}
\end{anhang}\newcommand{\dateiname}{L04401}\newcommand{\titel}{Arthur Schnitzler an Gustav Schwarzkopf, 18. 6. 1917}\newcommand{\editorInnen}{Herausgegeben von Jahnke, SelmaMüller, Martin Anton}\input{../tex-inputs/latex-leseansicht-abspann}
